% Presentation slides template for texforge (beamer, 16:9).
% Placeholders are filled in when the project is created; every % comment
% says what to replace. Text below is generic sample content.
% Tectonic reads UTF-8 directly, so no input encoding package is loaded.
\documentclass[aspectratio=169,11pt]{beamer}

% Document language from the {{language}} placeholder (default english).
\usepackage[{{language}}]{babel}
\usepackage[T1]{fontenc}
\usepackage{lmodern}
\usepackage{booktabs}
\usepackage{graphicx}
\usepackage{float}

% Built-in default theme; no external theme package (no metropolis).
\usetheme{default}

% One accent colour for frame titles, structure and itemize bullets.
\definecolor{acento}{HTML}{1F4E79}
\setbeamercolor{frametitle}{fg=acento}
\setbeamercolor{structure}{fg=acento}
\setbeamercolor{itemize item}{fg=acento}
\setbeamercolor{itemize subitem}{fg=acento}
\setbeamercolor{itemize subsubitem}{fg=acento}

% Navigation symbols off; footline keeps author, short title and frame/total.
\setbeamertemplate{navigation symbols}{}
\setbeamercolor{author in head/foot}{fg=black,bg=black!8}
\setbeamercolor{title in head/foot}{fg=black,bg=black!8}
\setbeamercolor{date in head/foot}{fg=black,bg=black!8}
\setbeamertemplate{footline}{%
  \leavevmode%
  \hbox{%
    \begin{beamercolorbox}[wd=.4\paperwidth,ht=2.5ex,dp=1ex,center]{author in head/foot}%
      \usebeamerfont{author in head/foot}\insertshortauthor
    \end{beamercolorbox}%
    \begin{beamercolorbox}[wd=.4\paperwidth,ht=2.5ex,dp=1ex,center]{title in head/foot}%
      \usebeamerfont{title in head/foot}\insertshorttitle
    \end{beamercolorbox}%
    \begin{beamercolorbox}[wd=.2\paperwidth,ht=2.5ex,dp=1ex,center]{date in head/foot}%
      \usebeamerfont{date in head/foot}\insertframenumber/\inserttotalframenumber
    \end{beamercolorbox}}%
  \vskip0pt%
}

% Replace the short forms with a short title and short author for the footline.
\title[Short title]{{{title}}}
\subtitle{{{subtitle}}}
\author[{{author}}]{{{author}}}
\institute[{{institution}}]{{{institution}}}
\date{{{date}}}

\begin{document}

  % Replace the title frame fields with your own; keep the placeholders.
  \begin{frame}
    \titlepage
  \end{frame}

  \begin{frame}{Outline}
    \tableofcontents
  \end{frame}

  \section{Motivation}

  % One point per line; \pause reveals them one at a time.
  \begin{frame}{Motivation}
    Why this work matters:
    \begin{itemize}
      \item The first point motivates the problem.
      \pause
      \item The second point names the gap.
      \pause
      \item The third point states the goal of the talk.
    \end{itemize}
  \end{frame}

  \section{Method}

  % Replace the diagram with your own; keep style= and the caption.
  \begin{frame}{Method}
    \begin{mermaid}[style=editorial, caption={Sample pipeline from data to result}, width=0.8\linewidth]
      flowchart LR
      data[Collect] --> clean[Clean] --> model[Train] --> report[Report]
    \end{mermaid}
  \end{frame}

  % Replace the listing with your own code (at most ten short lines).
  \begin{frame}[fragile]{Code}
    \begin{code}[lang=python]
def average(values):
    """Return the mean of a non-empty sequence."""
    return sum(values) / len(values)

scores = [0.7, 0.8, 0.9]
print(average(scores))
    \end{code}
  \end{frame}

  \section{Results}

  % Replace the rows with your own measurements.
  \begin{frame}{Results}
    \begin{center}
      \begin{tabular}{lcc}
        \toprule
        Run & Value & Error \\
        \midrule
        Baseline & 1.00 & 0.05 \\
        Variant A & 1.15 & 0.07 \\
        Variant B & \textbf{1.24} & 0.06 \\
        \bottomrule
      \end{tabular}
    \end{center}
    \begin{block}{Main result}
      The best value improves on the baseline by about a quarter.
    \end{block}
    \begin{alertblock}{Caveat}
      The sample numbers are illustrative; replace them with real data.
    \end{alertblock}
  \end{frame}

  % Two columns: text on the left, a diagram on the right.
  \begin{frame}{Two columns}
    \begin{columns}[T]
      \begin{column}{0.48\textwidth}
        Replace this paragraph with the explanation that belongs next to the
        diagram. Keep it short so both columns stay inside the frame.
      \end{column}
      \begin{column}{0.48\textwidth}
        \begin{graphviz}[style=monochrome, caption={Sample flow}]
          digraph G {
            rankdir=LR
            input -> process -> output
          }
        \end{graphviz}
      \end{column}
    \end{columns}
  \end{frame}

  \section{Conclusions}

  % One line per conclusion; keep the frame short.
  \begin{frame}{Conclusions}
    \begin{itemize}
      \item The method solves the stated problem.
      \item The result matches the sample measurements.
      \item The next step is to validate on real data.
    \end{itemize}
  \end{frame}

  % Replace the closing line with your own contact or question prompt.
  \begin{frame}{Thank you}
    \centering
    \Large Questions?
    \vspace{1em}

    \normalsize\texttt{author@example.org}
  \end{frame}

\end{document}
